\documentclass[11pt]{article}
\usepackage[margin=1in]{geometry}
\usepackage{amsmath,amssymb}
\usepackage[hidelinks]{hyperref}

\newcommand{\cat}{\mathrm{CAT}(0)}

\begin{document}

\begin{center}
{\Large Literature answer for arXiv:1807.01355}\\[0.4em]
{\large Modular semilattice orthoscheme complexes are \(\cat\)}
\end{center}

\paragraph{Status.}
\textbf{Literature already answered.}  Bacak's survey \emph{Old and new
challenges in Hadamard spaces} asks, in the section on submodular functions on
modular lattices, whether the orthoscheme complex \(K(\mathcal L)\) of every
finite-rank modular semilattice \(\mathcal L\) is a \(\cat\) space.  The survey
states that this is Chalopin--Chepoi--Hirai--Osajda Conjecture 7.3.

\paragraph{Source problem.}
After recalling that Chalopin--Chepoi--Hirai--Osajda proved the result for
finite-rank modular lattices, Bacak records the problem:
\[
\text{Let \(\mathcal L\) be a modular semilattice of finite rank. Is
\(K(\mathcal L)\) a \(\cat\) space?}
\]

\paragraph{Answering paper.}
Hirai's paper \emph{A Nonpositive Curvature Property of Modular Semilattices}
explicitly identifies the same conjecture.  In the introduction it says that
Chalopin et al. conjectured the CAT(0) property for modular semilattices and
adds ``see also [Problem 6.10]'' in Bacak's survey.  Its main theorem is:
\[
\text{If \(\mathcal L\) is a modular semilattice of finite rank, then
\(K(\mathcal L)\) is \(\cat\).}
\]
This is Theorem 1.2 in the arXiv version, and it is the same statement as the
source problem.

\paragraph{Scope.}
This packet clears only this orthoscheme-complex question from the source
survey.  It does not address the survey's other Hadamard-space problems, such
as weak topologies, compactness of convex hulls, stochastic proximal point
algorithms without local compactness, or asymptotic regularity of cyclic
projections.

\begin{thebibliography}{9}

\bibitem{Bacak2018}
Miroslav Bacak.
\newblock \emph{Old and new challenges in Hadamard spaces}.
\newblock arXiv:1807.01355, 2018.

\bibitem{CCHO}
Jeremie Chalopin, Victor Chepoi, Hiroshi Hirai, and Damian Osajda.
\newblock \emph{Weakly modular graphs and nonpositive curvature}.
\newblock arXiv:1409.3892; Memoirs of the AMS.

\bibitem{Hirai2019}
Hiroshi Hirai.
\newblock \emph{A Nonpositive Curvature Property of Modular Semilattices}.
\newblock arXiv:1905.01449; Geometriae Dedicata 214 (2021), 427--463.

\end{thebibliography}

\end{document}
